% 07-authorization.tex: verification is not authorization.
\section{Verification is not authorization: status, revocation and the offline trade}
\label{sec:authorization}

\motif{A credential's authenticity is not its current authorization.}

\analogy{A seal can be genuine and the traveller still unwelcome. The board by the gate, not the seal, says whose papers are honoured today, and a copy of the board carried along the road is trusted only until it goes stale.}

\subsection{Two results, never one}

The whole product is visible in one object. A relying party asks the issuer about a credential and
receives, for a revoked token:

\Needspace{9\baselineskip}
\begin{lstlisting}
POST /api/v1/verify  ->  {
  "authentic": true,                 // the ML-DSA signature is genuine
  "currently_authoritative": false,  // but this token has been revoked
  "status": "REVOKED",
  "usable": false,                   // authentic AND authoritative -> false
  "decision": "reject"
}
\end{lstlisting}

Authenticity is permanent and cacheable. Authorization is fresh and revocable. A relying party needs
both, and the API refuses to collapse them into one flag, because a signature stays genuine forever
while the authority under it can be withdrawn at any moment. \sImpl

The relying-party endpoint authenticates the caller as an organisation, not a person, by an OAuth2
client-credentials grant; it accepts a possession proof rather than a sequential identifier, so it
cannot be used to enumerate tokens; it returns a uniform verdict on a mismatch, so it is not an
existence oracle; and it keeps no record of who verified whom. \sImpl

\subsection{Fresh state, not a stale replica}

The verify endpoint is routed to a read replica for throughput, and the repository's history
records the mistake that split the two results permanently. Authenticity is a property of immutable
material, the signed value, the signature bytes and the stored key, so a lagging replica can neither
forge it nor deny it. Whether the token is usable \emph{now} is a current-authorization question, and
a revocation lands on the primary; a replica inside its staleness window could still show a
just-revoked token as active. The status that decides \code{usable} is pinned to
the primary, the response names its source and carries \code{as\_of} and a zero staleness bound, and
a caller may cache authenticity and must re-read authorization. \sImpl

\subsection{Revocation}

Revocation is a status transition through one sanctioned procedure that also writes to
\entity{RevocationList} in the same transaction, so the verifier-facing list can never diverge from
the token's state, and a trigger refuses any other path into \code{REVOKED}. The list carries a
canonical reason vocabulary; the audit trail carries the procedural detail and the co-signer when
the issuer-discretion ceiling required one. Revocation looks forward: events recorded while a
credential was valid are not invalidated, because history is not rewritten. \sImpl

\subsection{Authorization with the network off}

\figfloat{offline-status}{The offline trade, drawn as a timeline. The issuer signs a status
assertion with a window; the holder staples it; a relying party verifying offline inside the window
accepts. If the issuer revokes the credential inside the window, the last \code{ACTIVE} assertion
stays valid until it expires, and an offline verifier cannot know. The window is bounded by the
issuer's time-to-live (an hour by default) and by the \code{max\_window} the verifier imposes; a
high-risk relying party demands seconds or goes online. Who chooses the window, and who bears it: the
issuer sets the time-to-live and each relying party its \code{max\_window}; the harm of a long
window falls on whoever relies on a credential revoked inside it, and on the holder whose
revocation was meant to protect them. Polaris makes the window explicit and bounded; how long is
acceptable is a governance decision it does not make.}{fig:offline-status}

An offline verifier cannot ask the issuer. Polaris answers with a short-lived signed status
assertion: the issuer's signature, under the same key that signed the credential, over the token's
current status and a window. The holder fetches it by presenting the genuine credential, with no
bearer and no personal data, and staples it to a presentation. A verifier accepts offline only if
the credential is authentic, the assertion is authentic, bound to this credential, fresh on the
verifier's own clock, no wider than the window the verifier will accept, and \code{ACTIVE}. The rule
that a verifier never accepts a credential more than one window after its revocation, and that the
window is the one the verifier set rather than the one the issuer minted, is now a formal model
checked on every push with a configuration it must fail. \sImpl

The cost is stated as plainly as the benefit. Because verification touches no issuer, the issuer
never learns it happened: offline is more private than online. And because it touches no issuer, a
revoked credential's last assertion stays valid until it expires. Issuer non-observation, offline
availability and revocation freshness cannot all be maximised at once; a deployment picks two, and
the window is a policy number, not a law. One consequence is recorded on the scoreboard as a known
limitation rather than fixed: the issuer mints a one-hour window, the reference verifier's own
default bound is unset, and the ceiling the artifact itself carries is not read by the assertion
check, so the maximum offline lifetime of a status rests partly on the relying party's configuration.
A relying party that needs seconds sets seconds or goes online.

\subsection{Status across authorities, and through networks that are not trusted}

Two more signed objects carry status between authorities without a callback the issuer could log.
An epoch checkpoint is an authority's commitment to the latest point on its epoch chain, the epoch
number, its Merkle root and the prior epoch it extends; two checkpoints prove monotonicity, and two
different roots signed at one epoch number are cryptographic proof that the authority equivocated
about its own history. A revocation feed is the sorted set of revoked-credential leaves, each the
SHA3-256 of a token value, derivable only by the holder; a newer feed that drops a published
revocation is a caught rollback, because the underlying table is append-only. At federation scale an
aggregator mirrors many authorities' feeds and checkpoints, verbatim and still under each member's
own signature, into one short-lived bundle. The aggregator is untrusted for correctness: in full
control of the bundle, recomputing the commitment and re-signing the envelope, it still cannot forge
a member's status, because it cannot sign as the member. The bundle adds availability, not trust.
\sImpl

Signing a status artifact is also what lets an untrusted cache carry it, and a cache introduces the
one failure signing does not prevent, which is time. A cache directive is never a
constant: every windowed artifact's \code{max-age} is its own remaining life, computed from the
\code{expires\_at} the issuer signed, so a cache cannot outlive the window the issuer committed to;
every artifact that names one credential is \code{no-store}, because a shared cache holding one
would serve one holder's credential to another; and an expired or unparseable window is not cached at
all. Distribution at content-delivery scale is thereby specified and drilled; no deployment has it.
\sImplMark{} for the artifacts, drills and cache rules; \sFut{} for a distribution deployment.

The epoch cadence is the single number that sets three unrelated properties, and a design record now
says which: revocation freshness, since a revoked credential drops out of the \emph{next} epoch; the
size of the crowd, since the anonymity set is the epoch; and how often a relying party's
one-human-once ledger resets. Closing an epoch was once the expensive part, and it was expensive for
a bad reason: a fixed-depth tree padded to the national depth of 24 materialised sixteen million
leaves whatever the population was, so a root over a thousand members took 10.8 seconds and
2.9~GB. The padding is one repeated value hashed once per level, and a root over a
million members takes 2.1 seconds in tens of megabytes; the binding constraint on cadence is the
anonymity floor per context, not the compute. \sImpl

\begin{layercard}{Layer card: verification and status}
\qa{What it is}{The relying-party verify endpoint, the revocation list, the signed status assertion, epoch checkpoints, revocation feeds, the aggregate status bundle, and the cache rules that let an untrusted network carry them.}
\qa{Why it exists}{Authenticity is not authorization. A verifier needs to know not only that a claim was made but that it still stands, and sometimes needs to know it with the network off.}
\qa{What it trusts}{The issuer's signing key, for the assertion; its own clock, for freshness; the primary database, for the online verdict.}
\qa{What it refuses}{A replica's opinion of current status; an assertion wider than the verifier's ceiling; an aggregator's word for a member's status; a feed not bound to the issuer's key; a cache entry that would outlive the window the issuer signed.}
\qa{What it can see}{The status of one credential at one instant; the set of revoked leaves.}
\qa{What it cannot see}{A revocation that happened after the assertion it holds was signed; the active population, which no feed publishes.}
\qa{If it fails}{A stale verdict is bounded by the window, and the bound is the verifier's; a forked epoch chain or a rolled-back feed is non-repudiable evidence against the authority; a lying aggregator is caught at the member's signature.}
\qa{Now or later}{Now, for the artifacts and the rules. Distribution at national scale is later.}
\qa{Inside and outside}{Inside: the verifiers. Outside: the privacy boundary, which decides what a verifier may learn while it checks.}
\end{layercard}

\nextlayer{Status proves current authority, but proving it can reveal far more than the verifier
needed, and every verification is a fact about a person's day. The next layer bounds what anyone
learns, and this version reports what was measured about that bound rather than what was intended.}
